\section{R6: Contracts beyond finite decidable observations}\label{sec:r6}

\question{R6.Q1}{How should value construction and proof obligations be interleaved?}
\status{Use a dependency-aware obligation graph with cheap early propagation and delayed expensive proof search.}

The dichotomy ``prove immediately'' versus ``wait for a complete program'' is
too coarse. Some proof obligations can already determine data: solving
$?n=0$ by unification fixes a witness. Others are meaningful as symbolic local
obligations under an induction hypothesis. Expensive general automation on a
proposition whose relevant data holes are still unconstrained is usually
wasteful, but postponing all reasoning loses valuable type information.

For each obligation, compute its support: the value holes, type assignments,
local hypotheses, and instance choices on which it depends. Separate cheap
propagation rules---reflexivity, constructor injectivity, contradiction,
definitional simplification, and rigid index constraints---from open-ended
proof attempts. Run the former when their premises are available. Schedule
the latter after their support is stable, preferably at a branch boundary
rather than only at the root.

Refinement-type synthesis supplies a useful conceptual precedent for using
logical constraints while constructing terms \cite{synquid}. Without an SMT
solver, use a small finite vocabulary of predicates with Lean-checked
implication certificates. This is not automatically a complete liquid-type
inference procedure over arbitrary propositions; it is a restricted abstract
domain whose transfer rules are justified in Lean.

For length-preserving map, synthesize a recursive value with target
\[
 \Pi xs:\List A,\quad
 \{ys:\List B\mid \mathsf{length}(ys)=\mathsf{length}(xs)\}.
\]
The empty branch produces an empty list and a reflexive proof. The cons
branch takes the recursive pair, constructs the new head and tail, and uses
the recursive proof to discharge the successor-length equation. The proof is
local and follows the program, rather than being reconstructed from an opaque
completed recursor term.

This design needs a concrete implementation change. The inspected
\code{proofPortfolio} currently rejects goals whose type contains free locals,
not only unresolved metavariables. Local branch proofs therefore cannot simply
be routed through that function unchanged. Introduce a contextual proof
adapter using the goal's local context, with support-aware caching and
transactional rollback \cite{searchcode}. Completed proof terms may be cached
only after abstracting the locals they use.

\question{R6.Q2}{Can one recognize properties provable by program-aligned induction and arithmetic?}
\status{One can recognize a useful sufficient class by generating and classifying local verification conditions; a complete syntactic characterization is too strong.}

Let a program be a fold $q(xs)$ with initializer $z$ and step $s$. Supply an
invariant relation $I(xs,k)$. Its proof obligations are
\[
 I([],z),\qquad
 \forall a,xs,k,\ I(xs,k)\to I(a::xs,s(a,k)).
 \tag{8}\label{eq:inductionvc}
\]
An observation theorem then transfers $I(xs,q(xs))$ to the requested contract.
Indexed constructors give the analogous family of branch obligations.

\begin{theorem}[A certifiable program-aligned class]\label{thm:vc}
Fix a recursion scheme with its induction principle, a finite checked rewrite
system for constructor and program equations, and a sound decision procedure
for a theory $\mathcal D$. If an invariant and its generated base and step
obligations simplify to formulas accepted by $\mathcal D$, and the required
observation theorem is proved, then the synthesized program satisfies the
contract.
\end{theorem}
\begin{proof}
The checked rewrites preserve each obligation. The decision procedure supplies
proofs of all base and step obligations. Apply the recursion scheme's induction
principle to obtain the invariant at every input, then apply the observation
theorem. Every composition is represented by an ordinary proof term.
\end{proof}

This gives a positive recognizer: generate the obligations, perform the fixed
rewrites, and check whether the residual syntax belongs to the supported
theory. An initial theory can combine finite constructor case splits with
linear arithmetic over naturals and integers. Length accounting and many
simple measure bounds fit. Sorting additionally needs ordering and permutation
invariants; an accumulator-based reverse proof often needs an associativity
lemma and a generalized accumulator statement. Their failure to fall into the
initial class does not establish that they need fundamentally different
induction.

Membership in a narrowly defined residual theory can be decided syntactically.
Whether \emph{some} invariant, generalization, rewrite sequence, or library
lemma would make an arbitrary Lean property fit is a new search problem.
The engine should report ``outside this certified induction tier,'' not
``not inductively provable.'' Store successful invariants as reusable proof
patterns, with their exact conditions, rather than expanding the default
simplifier indiscriminately.

\question{R6.Q3}{What induction automation is usable in Lean now?}
\status{Core induction and program-aligned proof construction are immediately usable; broader automated induction remains a portfolio component rather than a universal service.}

For the pinned toolchain, \code{MVarId.induction} is an available meta-level
library operation, while \code{simp}, \code{omega}, and \code{grind} can serve
as branch dischargers. Lean's library-search implementation itself exposes
an example of invoking \code{Grind.main} on a goal with default parameters.
This is a concrete integration path, not merely a suggestion to invoke an
interactive tactic from a string \cite{leaninduction,librarysearch}.

The current official manual, consulted as version 4.35.0-rc2, also documents
\code{fun_induction}: it uses generated functional induction principles for
non-mutual structurally or well-founded recursive definitions, supports
\code{generalizing}, and can unfold the function call while setting up the
proof. \code{fun_cases} provides the corresponding case analysis
\cite{leantactics}. This is particularly relevant after publishing a candidate
through the equation compiler. Check availability against Leant2's pinned
4.34.0 build rather than silently treating the moving manual as that build.

Keep induction selection separate from branch automation. A program-guided
adapter knows the recursive argument and the equations that generated the
body; it can choose the corresponding induction principle and present the
right local hypotheses. Calling \code{grind} on a closed theorem without that
structure should not be assumed to discover the required induction and
auxiliary invariant. Conversely, applying induction to every inductive local
reintroduces the search explosion already observed in the proposal.

Canonical supplies an additional dependent inhabitation and recursion baseline,
Aesop a rule-based proof-search framework, and LeanHammer a learned
premise-selection plus proof-reconstruction approach
\cite{canonical,aesop,leanhammer}. These have different dependencies and trust
boundaries. A core-only Leant2 default should use a small internal adapter and
allow optional integrations through the same proof-service interface. No
current benchmark for those tools establishes ten-second success on Leant2's
quantified synthesis contracts.

A useful service reply is a proof, a checked counterexample, an unsupported
fragment, or resource exhaustion. Give each invocation a heartbeat quota and
a cancellation boundary; an external process may additionally need an
operating-system limit. A wall-clock deadline cannot necessarily interrupt
all code at an arbitrary instruction. Record timeout separately from logical
failure. Rippling-inspired generalization and lemma discovery should initially
produce candidates for this service, not silently extend its trusted calculus.

\question{R6.Q4}{Is there a principled scheduler for proof debt?}
\status{Yes as a constrained decision policy with explicit costs and fairness; no simple bandit index is universally optimal for these dependent obligations.}

For an obligation $o$, estimate its remaining cost $\widehat t(o)$, chance of
usefully resolving $\widehat p(o)$, and downstream benefit $\widehat b(o)$.
A provisional priority is
\[
 \frac{\widehat p(o)\,\widehat b(o)}
      {\widehat t(o)+\widehat s(o)+\varepsilon},
 \tag{9}\label{eq:debt}
\]
where $\widehat s$ models setup or switching cost. Benefit should include
shared consequences: proving one lemma may discharge several candidates,
while rejecting an early local invariant may eliminate a large subtree.
Equation~\eqref{eq:debt} is a heuristic whose assumptions are visible, not an
optimality theorem.

These jobs are not stationary independent arms. Filling a data hole changes
several proof problems; an arithmetic lemma can make all of them easy; a
counterexample changes the synthesis frontier. An exact policy could be
formulated over a finite state model with known transition probabilities and
costs, but learning that model is itself expensive. Classic restart results
likewise do not automatically cover this dependency structure \cite{luby}.

Use a deterministic fair scheduler first: counterexamples before expensive
proofs, local obligations before global reconstruction when they are ready,
small proof quanta before larger ones, and a reserved share for synthesis.
Attach each proof job to the support fingerprint on which it was created.
Suspend rather than restart jobs when the solver permits it. Learn estimates
only at logical epoch boundaries, and evaluate the learned policy against
fixed round-robin and fixed-tier baselines. Measure total certified solutions
and time to rejection, not the number of tactic calls that report success.
